%% kapa: exam=2 round=37 year=2026 period=3 subject=compensation_law subject_ko=감정평가 및 보상법규 seq=3 number=3 points=20
%% source: Q-Net 감정평가사 기출문제 (cst00309) artlSeq=5268128 file=raw/1008/5268128/26년 제37회 감정평가사 2차 (3교시_감정평가 및 보상 법규).pdf
%% conversion: kordoc 4.24.51 -> tools/gen2.py (auto); review=proofread-gpt-terra; flags=none
\begin{longquestion}{【문제 3】}
건축주 A는 X도 B시에 대지면적 900m\textsuperscript{2}인 주거용 건축물(이하 ‘이 사건 건축물’이라 한다)을 신축하여 소유권을 취득하였고 일반취득세율을 적용하여 취득세를 신고납부하였다. 당시 「지방세법 시행령」에 따르면 대지면적 662m\textsuperscript{2}를 초과하고 취득 당시의 시가표준액이 9억원을 초과하는 경우 취득세 중과세 대상인 ‘고급주택’에 해당하였다. 이 사건 건축물은 아직 개별주택가격이 공시되지 않은 상태였기에 B시장은 국토교통부장관이 제공한 ‘주택가격비준표’와 ‘개별주택가격 조사ㆍ산정지침’을 근거로 이 사건 건축물의 시가표준액을 9억6천만원으로 산정하였다. 이에 따라 B시장은 A에게 이 사건 건축물이 ‘고급주택’ 기준을 초과한다는 이유로 중과세율을 적용하여 취득세 및 신고ㆍ납부의무 위반에 따른 가산세 부과처분예고를 하였다. \par
A는 ‘주택가격비준표’ 및 ‘개별주택 조사ㆍ산정지침’은 대외적 구속력이 없으므로 이 사건 건축물은 고급주택이 아니라고 주장한다. ‘주택가격비준표’ 및 ‘개별주택 조사ㆍ산정지침’의 법적 근거와 법적 성격에 대해서 서술하시오. (20점) \par
\begin{framedtext}
참조 조문 \par
<지방세법> \par
제4조(부동산 등의 시가표준액) ① 이 법에서 적용하는 토지 및 주택에 대한 시가표준액은 「부동산 가격공시에 관한 법률」에 따라 공시된 가액(價額)으로 한다. 다만, 개별공시지가 또는 개별주택가격이 공시되지 아니한 경우에는 특별자치시장ㆍ특별자치도지사ㆍ시장ㆍ군수 또는 구청장(자치구의 구청장을 말한다. 이하 같다)이 같은 법에 따라 국토교통부장관이 제공한 토지가격비준표 또는 주택가격비준표를 사용하여 산정한 가액으로 하고, 공동주택가격이 공시되지 아니한 경우에는 대통령령으로 정하는 기준에 따라 특별자치시장ㆍ특별자치도지사ㆍ시장ㆍ군수 또는 구청장이 산정한 가액으로 한다. \par
② 제1항 외의 건축물(새로 건축하여 건축 당시 개별주택가격 또는 공동주택가격이 공시되지 아니한 주택으로서 토지부분을 제외한 건축물을 포함한다), 선박, 항공기 및 그 밖의 과세대상에 대한 시가 표준액은 거래가격, 수입가격, 신축ㆍ건조ㆍ제조가격 등을 고려하여 정한 기준가격에 종류, 구조, 용도, 경과연수 등 과세대상별 특성을 고려하여 대통령령으로 정하는 기준에 따라 지방자치단체의 장이 결정한 가액으로 한다. \par
③ ∼ ④ (생 략) \par
\end{framedtext}
\end{longquestion}
